\begin{figure}[H]
\centering
\includegraphics[width=.78\linewidth]{figures/lens-shaped-region-source.png}
\caption{A lens-shaped region}\label{fig:holmgren-lens-region}
\end{figure}

\subsection{Holmgren's Uniqueness Theorem}\label{sec:holmgren-uniqueness}

The theorem in the previous section shows existence and uniqueness of solutions for a noncharacteristic initial-value problem. However, uniqueness was only guaranteed within the class of analytic functions; the existence of other, nonanalytic solutions was not ruled out. Holmgren's theorem shows that this cannot happen for linear equations; we shall prove uniqueness assuming only that the solution is smooth enough so that all derivatives appearing in the partial differential equation are continuous (using the concept of "generalized" solutions, defined later in this book, this assumption can be relaxed further). The proof of uniqueness is achieved by proving existence of solutions for an "adjoint" system of differential equations. To obtain this existence we shall use the Cauchy-Kovalevskaya theorem; this requires us to assume analytic coefficients in the equations. If, however, we had an existence theory which works without analyticity of the coefficients, this assumption would be unnecessary.

\subsubsection{An Outline of the Main Idea}\label{subsec:holmgren-main-idea}

Consider a system of linear equations
\begin{equation}\label{eq:holmgren-linear-first-order-system}
a_{ij}^k(\bm{x})\dfrac{\partial u_j}{\partial x_k}+b_{ij}(\bm{x})u_j=0,\qquad i=1,\ldots,N.
\end{equation}
Let $\bm{u}=(u_1,\ldots,u_N)$ be a solution in a "lens-shaped" domain $\Omega\subset\mathbb{R}^n$ bounded by two surfaces $S$ and $Z$. Assume that $\bm{u}=0$ on $Z$ and that $S$ is noncharacteristic and analytic. We also assume that the coefficients in \eqref{eq:holmgren-linear-first-order-system} are analytic.

Let $v_i$, $i=1,\ldots,N$ be arbitrary functions in $C^1(\overline\Omega)$. We multiply the $i$th equation of \eqref{eq:holmgren-linear-first-order-system} by $v_i$, sum over $i$, and integrate over $\Omega$. This yields
\begin{equation}\label{eq:holmgren-integration-by-parts}
\begin{aligned}
0&=\int_\Omega v_i(\bm{x})a_{ij}^k(\bm{x})\dfrac{\partial u_j}{\partial x_k}(\bm{x})+
v_i(\bm{x})b_{ij}(\bm{x})u_j(\bm{x})\,\mathrm{d}\bm{x}\\
&=\int_\Omega-\dfrac{\partial}{\partial x_k}[v_i(\bm{x})a_{ij}^k(\bm{x})]u_j(\bm{x})
+v_i(\bm{x})b_{ij}(\bm{x})u_j(\bm{x})\,\mathrm{d}\bm{x}\\
&\quad+\int_{\partial\Omega}a_{ij}^k(\bm{x})v_i(\bm{x})u_j(\bm{x})n_k\,\mathrm{d}S,
\end{aligned}
\end{equation}
where $\bm{n}$ is the outer normal to $\partial\Omega$.

Assume now that $\bm{v}$ satisfies the "adjoint" system of PDEs,
\begin{equation}\label{eq:holmgren-adjoint-system}
-\dfrac{\partial}{\partial x_k}(a_{ij}^kv_i)+b_{ij}v_i=0,\qquad j=1,\ldots,N,
\end{equation}
with initial conditions
\begin{equation}\label{eq:holmgren-adjoint-cauchy-data}
v_i=f_i
\end{equation}
on $S$. Then \eqref{eq:holmgren-integration-by-parts} reduces to
\begin{equation}\label{eq:holmgren-boundary-pairing}
0=\int_Sa_{ij}^k(\bm{x})f_i(\bm{x})u_j(\bm{x})n_k\,\mathrm{d}S.
\end{equation}


If this holds for arbitrary continuous functions $f_i$ on $S$, then we conclude that $a_{ij}^ku_jn_k=0$ on $S$, and since $\det a_{ij}^kn_k\ne0$ ($S$ is noncharacteristic), we conclude that $\bm{u}=0$ on $S$.

The Cauchy-Kovalevskaya theorem guarantees that \eqref{eq:holmgren-adjoint-system} has a solution in a neighborhood of $S$ if the $f_i$ are analytic. Unfortunately, we can in general not claim that this neighborhood includes all of $\Omega$. If it did, we would obtain \eqref{eq:holmgren-boundary-pairing} for analytic $\bm{f}$. The Weierstrass approximation theorem states that any continuous function on a compact subset of $\mathbb{R}^n$ can be approximated uniformly by polynomials. Therefore, if \eqref{eq:holmgren-boundary-pairing} holds for $\bm{f}$ whose components are polynomials, it also holds for continuous $\bm{f}$.

\subsubsection{Statement and Proof of the Theorem}\label{subsec:holmgren-statement-proof}

In order to overcome the difficulty that we cannot guarantee a solution of \eqref{eq:holmgren-adjoint-system} throughout all of $\Omega$, we shall replace the surface $S$ by a one-parameter family of surfaces $S_\lambda$ and then take "small steps" in $\lambda$. More precisely, we shall presume the following situation.

Let $D$ be a bounded domain in $\mathbb{R}^n$, such that the coefficients of \eqref{eq:holmgren-linear-first-order-system} are analytic on $D$. Let $Z=D\cap\{x_n=0\}$ and assume that $Z$ is nonempty and noncharacteristic. Let $\Phi(\bm{x})$ be an analytic function defined on $D$ such that $\nabla\Phi\ne0$ and let $S_\lambda=\{\Phi=\lambda\}\cap\{x_n\geq0\}$. We assume there are real numbers $a$ and $b$, $a<b$, such that the following hold:

\begin{enumerate}
\item The set $\bigcup_{\lambda\in[a,b]}S_\lambda$ is compact.
\item $S_a$ consists of a single point located on $Z$.
\item For $a<\lambda<b$, $S_\lambda$ is a regular surface intersecting $Z$ transversally (the intersection of two surfaces is called transverse if their normals are not collinear). The intersection of $S_\lambda$ and $Z$ is then a regular (analytic) $(n-2)$-dimensional surface. Moreover, we assume that $S_\lambda$ is noncharacteristic.
\end{enumerate}

\begin{figure}[H]
\centering
\includegraphics[width=.78\linewidth]{figures/holmgren-surface-family-source.png}
\caption{}\label{fig:holmgren-surface-family}
\end{figure}

We shall establish the following result:

\begin{Theorem}
\label{thm:holmgren-uniqueness}
Let $\Omega=\{\bm{x}\in D\mid x_n>0,\ a<\Phi(\bm{x})<b\}$. Let $\bm{u}\in C^1(\overline\Omega)$ be a solution of \eqref{eq:holmgren-linear-first-order-system} such that $\bm{u}=0$ on $\partial\Omega\cap\overline Z$. Then $\bm{u}=0$ in $\overline\Omega$.
\end{Theorem}

\begin{Proof}Let $\Lambda=\{\lambda\in[a,b]\mid\bm{u}=0\text{ on }S_\lambda\}$. We know that $a\in\Lambda$, and it follows from the continuity of $\bm{u}$ that $\Lambda$ is closed. We shall show that $\Lambda$ is also open in $[a,b]$. This implies $\Lambda=[a,b]$ and hence the theorem.

We note that $\overline\Omega$ is compact, and hence there is $M,\rho$, independent of $\bm{x}$, such that $a_{ij}^k$, $b_{ij}$ and $\Phi$ are in $C_{M,\rho}(\bm{x})$ for every $\bm{x}\in\overline\Omega$. Consequently, if Cauchy data of class $C_{M,\rho}$ are prescribed on $S_\mu$, a solution of \eqref{eq:holmgren-adjoint-system} exists in an $\epsilon$-neighborhood of $S_\mu$, with $\epsilon$ independent of $\mu\in(a,b]$. We note that any polynomial lies in some $C_{M,\rho}$, where we can choose $\rho$ as large as we wish, at the expense of making $M$ large. However, \eqref{eq:holmgren-adjoint-system} is linear, and hence the domain on which the solution exists does not change if the Cauchy data are multiplied by a constant factor. Hence the class of Cauchy data for which solutions to \eqref{eq:holmgren-adjoint-system} exist in an $\epsilon$-neighborhood of $S_\mu$ includes all polynomials.

We claim that, for given $\lambda\in[a,b]$ and $\epsilon>0$, there is a $\delta>0$ such that $S_\lambda$ is contained in the $\epsilon$-neighborhood of $S_\mu$ whenever $\mu\in[a,b]$ and $|\mu-\lambda|<\delta$. To see this, we first note that, in the neighborhood of any point $\bm{x}\in S_\lambda$, the equation $\Phi(\bm{x})=\mu$ can be solved for one of the coordinates $x_i=x_i(x_1,\ldots,x_{i-1},x_{i+1},\ldots,x_n,\mu)$, and if $\bm{x}\in Z$ and $\lambda\ne a$ we can choose $i\ne n$. If $\lambda=a$, we have to choose $i=n$, and $x_n$ is an increasing function of $\mu$. In all cases, an immediate consequence is that, if $\delta(\bm{x})$ is chosen sufficiently small, then for every $\mu\in[a,b]$ with $|\mu-\lambda|<\delta(\bm{x})$ there is a point $\bm{y}\in S_\mu$ with $|\bm{y}-\bm{x}|<\epsilon/2$. Since $S_\lambda$ is compact, there is a finite number of points $\bm{x}_k$, $k=1,\ldots,K$, such that $S_\lambda$ is covered by the balls centered at $\bm{x}_k$ with radius $\epsilon/2$. The claim then follows with $\delta=\min\delta(\bm{x}_k)$.

Assume now that $\lambda\in\Lambda$ and $\mu\in(a,b]$ with $|\mu-\lambda|<\delta$, where $\delta$ is as above. We can then apply the argument explained in the previous section to the domain bounded by $S_\lambda$, $S_\mu$ and $Z$. We thus reach the conclusion that $\bm{u}=0$ on $S_\mu$, and hence $\mu\in\Lambda$.\end{Proof}

\begin{Example}
\label{ex:wave-cauchy-domain-of-determinacy}
Consider the wave equation in two dimensions
\begin{equation}\label{eq:wave-determinacy-example}
u_{yy}-u_{xx}=0,
\end{equation}
with Cauchy data prescribed for $y=0$, $-1<x<1$. Let $\Phi(x,y)=(x-y+1)(x+y-1)$, and let $D=(-1,1)\times(-1,1)$. Then $S_\lambda$, $-1\leq\lambda<0$ is the arc of the hyperbola $(x-y+1)(x+y-1)=\lambda$ that lies within the triangle with corners $(-1,0)$, $(1,0)$ and $(0,1)$. It is easy to show that all the hypotheses of the theorem are satisfied with $a=-1$ and any $b\in(-1,0)$. Since the $S_\lambda$ fill the interior of the triangle, $u$ is determined within the whole triangle by its prescribed Cauchy data. In general, if $u$ is determined in $\Omega$ by its Cauchy data on $Z$, we call $\Omega$ a \emph{domain of determinacy} for $Z$.
\end{Example}

\subsubsection{The Weierstrass Approximation Theorem}\label{subsec:weierstrass-approximation}

In the above proof, we have used the following theorem, known as the Weierstrass approximation theorem.

\begin{Theorem}
\label{thm:weierstrass-approximation}
Let $S$ be a compact subset of $\mathbb{R}^n$ and let $f$ be a continuous function $S\to\mathbb{R}$. Then there is a sequence of polynomials $p_m(\bm{x})$ such that $p_m(\bm{x})\to f(\bm{x})$ uniformly on $S$.
\end{Theorem}

\begin{Proof}Let $a_i$, $b_i$, $i=1,\ldots,n$ be such that $S\subset P:=\prod_{i=1}^n(a_i,b_i)$. By the extension theorem of Tietze-Urysohn, $f$ can be extended to a continuous function on all of $\mathbb{R}^n$, which agrees with the given function on $S$ and vanishes outside of $P$. (Most good topology texts give a proof of this result.) We shall again call the extended function $f$ and show that $f$ can be approximated by polynomials uniformly on $P$. Without loss of generality, we may assume that $0<a_i<b_i<1$.

We first consider the one-dimensional case, i.e., let $f\in C(\mathbb{R})$ be such that $f=0$ outside the interval $(a,b)$, where $0<a<b<1$. We shall construct a sequence $p_m$ of polynomials such that $p_m\to f$ uniformly on $[a,b]$.

First, we define
\begin{equation}\label{eq:weierstrass-kernel-integrals}
I_m=\int_0^1(1-v^2)^m\,\mathrm{d}v,\qquad I_{m,\delta}=\int_\delta^1(1-v^2)^m\,\mathrm{d}v.
\end{equation}
From the elementary inequalities
\begin{equation}\label{eq:weierstrass-kernel-integral-bounds}
I_m>\int_0^1(1-v)^m\,\mathrm{d}v=\dfrac1{m+1},\qquad I_{m,\delta}<(1-\delta^2)^m,
\end{equation}
we conclude that, for any positive $\delta$,
\begin{equation}\label{eq:weierstrass-kernel-tail-limit}
\lim_{m\to\infty}\dfrac{I_{m,\delta}}{I_m}=0.
\end{equation}


We now choose $\alpha$ and $\beta$ such that $0<\alpha<a<b<\beta<1$. We set
\begin{equation}\label{eq:weierstrass-one-dimensional-polynomial}
p_m(x)=\dfrac{\int_\alpha^\beta f(y)[1-(y-x)^2]^m\,\mathrm{d}y}{\int_{-1}^1(1-y^2)^m\,\mathrm{d}y}.
\end{equation}
Obviously, $p_m$ is a polynomial of degree $2m$. We shall now show that $p_m$ converges to $f$, uniformly on $[a,b]$. In the numerator, we set $y=v+x$, and obtain
\begin{equation}\label{eq:weierstrass-three-integral-split}
\begin{aligned}
\int_{\alpha-x}^{\beta-x}f(v+x)(1-v^2)^m\,\mathrm{d}v
&=\int_{\alpha-x}^{-\delta}\cdots+\int_{-\delta}^{\delta}\cdots+\int_\delta^{\beta-x}\cdots\\
&=I_1+I_2+I_3.
\end{aligned}
\end{equation}


Evidently, $|I_1|$ and $|I_3|$ are bounded by $MI_{m,\delta}$, where $M$ is the maximum of $|f|$. Let now $\epsilon>0$ be given and choose $\delta$ such that $|f(v+x)-f(x)|\leq\epsilon$ for $|v|\leq\delta$. Then we compute
\begin{equation}\label{eq:weierstrass-central-integral}
\begin{aligned}
I_2&=f(x)\int_{-\delta}^{\delta}(1-v^2)^m\,\mathrm{d}v\\
&\quad+\int_{-\delta}^{\delta}(f(x+v)-f(x))(1-v^2)^m\,\mathrm{d}v.
\end{aligned}
\end{equation}
The first term on the right equals $2f(x)(I_m-I_{m,\delta})$, and the second term can be estimated by $2\epsilon I_m$. By combining the estimates obtained, we find
\begin{equation}\label{eq:weierstrass-uniform-error-estimate}
|p_m(x)-f(x)|\leq2M\dfrac{I_{m,\delta}}{I_m}+\epsilon.
\end{equation}


We can make this quantity as small as we wish by first choosing $\epsilon$ small, then choosing $\delta$ accordingly and then choosing $m$ sufficiently large.

In several space dimensions, we proceed in an analogous fashion with
\begin{equation}\label{eq:weierstrass-multivariable-polynomial}
p_m(\bm{x})=\dfrac{\int_{\alpha_1}^{\beta_1}\cdots\int_{\alpha_n}^{\beta_n}f(\bm{y})[1-(y_1-x_1)^2]^m\cdots[1-(y_n-x_n)^2]^m\,\mathrm{d}y_n\cdots\mathrm{d}y_1}{\left[\int_{-1}^1(1-\xi^2)^m\,\mathrm{d}\xi\right]^n}.
\end{equation}
\end{Proof}

\paragraph*{Problems}

\begin{enumerate}[label=\thesection.\arabic*.,ref=\thesection.\arabic*,start=18,series=book-problems]

\item\label{ex:laplace-cauchy-uniqueness} Let $u$ be a solution of Laplace's equation in $\mathbb{R}^2$ and assume that $u=u_y=0$ for $y=0$ and $-1<x<1$. Show that $u=0$ everywhere.

\item\label{ex:wave-determinacy-boundary} Consider the wave equation $u_{yy}-u_{xx}=0$ with Cauchy data on $(-1,1)\times\{0\}$. Show that no domain of determinacy extends beyond the square with the corners $(-1,0)$, $(0,1)$, $(1,0)$ and $(0,-1)$.

\item\label{ex:constant-coefficient-cauchy-uniqueness} Consider a system of linear homogeneous first-order PDEs with constant coefficients such that the planes $x_n=\mathrm{const.}$ are noncharacteristic. Show that any solution which vanishes on $x_n=0$ vanishes everywhere.

\item\label{ex:nonunique-zero-initial-heat-solution} Verify that
\begin{equation}\label{eq:flat-in-time-heat-series}
u(x,t)=\sum_{m=0}^{\infty}\dfrac{x^{2m}}{(2m)!}\dfrac{\mathrm{d}^m}{\mathrm{d}t^m}\exp\left(\dfrac{-1}{t^2}\right)
\end{equation}
is a solution of $u_t=u_{xx}$ with initial condition $u(x,0)=0$. Why does Holmgren's theorem not apply? Hint: Use Cauchy's formula to estimate derivatives of $\exp(-1/t^2)$. For the contour, choose a circle centered at $t$ with radius $t/2$.

\item\label{ex:polynomial-approximation-on-bounded-sets} Let $f$ be a continuous function on $\mathbb{R}^n$. Show that there is a sequence of polynomials $p_m$ such that $p_m\to f$ on all of $\mathbb{R}^n$, uniformly on any bounded set. Give an example showing that in general $p_m$ cannot converge to $f$ uniformly on all of $\mathbb{R}^n$.

\end{enumerate}
